\documentclass[a4paper,12pt]{article}

\usepackage[french]{babel}
\usepackage[T1]{fontenc} 
\usepackage[latin1]{inputenc}
%\usepackage{amsmath}

\author{R�mi Synave, Stefka Gueorguieva, Pascal Desbarats}
\title{Manuel de la biblioth�que TOPOLOGY}
\date{}

\begin{document}
\maketitle

\section{Utilisation}
Des fonctions permettant des op�rations topologiques sur les maillages sont impl�ment�es dans cette biblioth�que.\\

\section{Fonctions}

\textbullet void vf\_model\_star(int type, vf\_model m, int *faces, int nbfaces, int **list, int *size, int depth)\\
Recherche du voisinage � partir d'une liste d'index de faces, le tableau \textit{faces} et en sp�cifiant la profondeur \textit{depth}.\\
Il est possible de demander le voisinage par ar�tes (\textit{type==BYEDGE}) ou par sommet (\textit{type==BYVERTEX}). La liste des faces appartenant au voisinage est stock�e dans \textit{list}.\\


\textbullet int vf\_model\_nb\_hole(vf\_model *m)\\
Retourne le nombre de trou dans le maillage.\\
Le nombre de trou est �gale au nombre de bords ferm�s.\\


\textbullet int vf\_model\_nb\_connected\_part(vf\_model *m, int **list)\\
Retourne le nombre de parties connexes contenu dans le maillage. Le param�tre \textit{list} contient le num�ro de la partie de la face du m�me index.\\

\end{document}
